\section{Projection constraints from two frames}
\label{sec:projections}

We retain the stationary data of
Proposition~\ref{prop:stationary-data} and the entropy rates of
Section~\ref{sec:entropy}. Thus
\[
 0<\tau<1,\qquad \ell=1-\tau,\qquad
 s_0=1-\beta>0,\qquad
 (k_X,k_Y,k_U,k_V)=(0,\tau,\ell,1).
\]
The angular exponent supplied by Lemma~\ref{lem:angular-frostman}
is denoted by $\gamma>0$.
All conditional probabilities below refer to the indexed event
probability, with its inherited weights.

\begin{proposition}[Projection constraints]
\label{prop:projection-rates}
At almost every interior point $(c,a)$ where the full-cell thresholds
lie in the permitted range, the following assertions hold.
In every case write $a_1=j_X$ and $v_1=j_V$.
If $\tau\ne\ell$, write $y_1=j_Y$ and $u_1=j_U$,
where the conditional rates are ordered by increasing speed $k_i$.
Each belongs to $[0,1]$, and
\begin{equation}
\begin{split}
 a_1&\ge y_1,\qquad u_1\ge v_1,\\
 v_1>0&\ \Longrightarrow\ y_1\ge\min\{1,s_0+v_1\},\\
 a_1<1&\ \Longrightarrow\ u_1\le(a_1-s_0)_+.
\end{split}
\label{pr:distinct-constraints}
\end{equation}
If $\tau=\ell$, put $m_1=j_{\{Y,U\}}$.
Then $a_1,v_1\in[0,1]$, $m_1\in[0,2]$, and
\begin{equation}
 a_1<1\quad\Longrightarrow\quad
 m_1\le1+(a_1-s_0)_+,
 \qquad v_1\le(a_1-s_0)_+.
\label{pr:tied-constraints}
\end{equation}
These conclusions apply on both signs of the offset whenever those
offsets are in the permitted range.
\end{proposition}

The geometric tests compare two frames of the same matrix. Write
$\Delta t$ and $\Delta\vartheta$ for the changes in time and normal
parameter between them. A time change adds velocity to position, and
a normal change adds the transverse component to the normal component. Thus, for example,
$Y'=Y+\Delta t\,V$ projects the pair $(Y,V)$ onto a line. If the
second-frame coordinate occupied too few bins, the planar projection
bound would force the frame label to reveal a positive amount of matrix
information. An averaged information estimate will rule this out.

We first put these tests in the units of a short scale interval. The
coincidence of the two middle speeds requires one additional projection,
whose slope is $\Delta t\Delta\vartheta$. We state the pinning input for
that test before choosing the pair laws. Its proof, together with the
planar projection argument used here, is given in
Section~\ref{sec:projection-proofs}.

\subsection{The local matrix experiment}

Fix an interior $c$ for which the entropy conclusions hold at almost
every offset, and work on a compact interval of admissible offsets.
Our base cell will stagger the four coordinate-refinement intervals:
coordinates with different speeds will be refined on disjoint depth
intervals. This makes the ordered conditional rates of
Section~\ref{sec:entropy} the dimensions of the source pairs used below.
We shall let $b\downarrow0$ and write
\[
 E_0=e_{c+b},\qquad D=N^b.
\]
In the $E_0$-frame, let $J_0$ be the full matrix cell at thresholds
\begin{equation}
 r_i^0=\rho+(w-k_i)c+a+(1-k_i)b,
 \qquad i\in\{X,Y,U,V\}.
\label{pr:rebase-thresholds}
\end{equation}
Equivalently, this is the full cell based at $c+b$ and offset
$a+(1-w)b$. The cell partitions are nested as $a$ increases.

Choose a fixed rational $\zeta_{\max}>0$ smaller than every nonzero
difference $|k_i-k_j|$. Choose a fixed $C_0>1$ sufficiently large,
and put
\begin{equation}
 L=e_{c+C_0b}.
\label{pr:fine-label}
\end{equation}
For small $b$ all labels remain in the original scale interval.
Write $\mathfrak s_i=N^{-r_i^0}$ for the four physical cell widths.
Subtract a fixed matrix center of $J_0$ and express each form in its
corresponding units $\mathfrak s_i$.
Here $X,Y$ are normal and transverse position, and $U,V$ are the
corresponding velocity coordinates, as in
\eqref{eq:stationary-matrix-forms}. The centered coordinates stay in a
fixed bounded set in every frame whose label lies within $E_0$.
A time change adds a multiple of the velocity row of
$\left(\begin{smallmatrix}X&Y\\ U&V\end{smallmatrix}\right)$
to its position row; a normal change adds a multiple of its transverse
column to its normal column. Their composition gives the mixed term
in the following formula.

\begin{lemma}[Change of frame in relative units]
\label{lem:pr-relative-frame}
If two frame labels in $E_0$ have differences $\Delta t$ and
$\Delta\vartheta$, measured respectively in units $R_{c+b}$ and
$\Theta_{c+b}$, their centered relative coordinates are related by
\begin{equation}
 (X,Y,U,V)\longmapsto
 \bigl(X+\Delta t\,U+\Delta\vartheta\,Y
                 +\Delta t\Delta\vartheta\,V,
       Y+\Delta t\,V,
       U+\Delta\vartheta\,V,V\bigr).
\label{pr:relative-frame}
\end{equation}
All coefficients are bounded. Reading the frame with the label
$L$ introduces relative errors
$O(D^{-\min(\ell,\tau)(C_0-1)})$ on bounded coordinates.
Thus $C_0$ can be fixed so that these errors are smaller than
$D^{-\zeta_{\max}}$ by a fixed positive power.
\end{lemma}

\begin{proof}
The widths in \eqref{pr:rebase-thresholds} satisfy
\begin{equation}
 \frac{\mathfrak s_X}{\mathfrak s_U}=R_{c+b},\quad
 \frac{\mathfrak s_X}{\mathfrak s_Y}=\Theta_{c+b},\quad
 \frac{\mathfrak s_Y}{\mathfrak s_V}=R_{c+b},\quad
 \frac{\mathfrak s_U}{\mathfrak s_V}=\Theta_{c+b},\quad
 \frac{\mathfrak s_X}{\mathfrak s_V}=R_{c+b}\Theta_{c+b}.
\label{pr:width-ratios}
\end{equation}
The last equality uses $\tau+\ell=1$.
Substituting these ratios in the exact change of the four linear
forms gives \eqref{pr:relative-frame}.
Using uncentered forms adds only known translations.
The time and angular errors of $L$ relative to $E_0$ are respectively
$D^{-\ell(C_0-1)}$ and $D^{-\tau(C_0-1)}$.
The formula, on a bounded set of coordinates and coefficients,
gives the asserted error bound.
Actual dyadic widths may be used as the units, or bounded rounding
factors may be retained throughout.
\end{proof}

For $\mathbf u=(u_X,u_Y,u_U,u_V)$ with
$0\le u_i\le\zeta_{\max}$, let $\mathcal O_{\mathbf u}$ observe
axis $i$ in the $L$-frame to relative depth $u_i$, that is, to
threshold $r_i^0+bu_i$.
An observation at depth zero carries at most bounded additional
entropy after $E_0,J_0,L$ are known. Figure~\ref{fig:entropy-layers}
shows the refinement intervals when all four depths equal a fixed
$\zeta\in(0,\zeta_{\max}]$.

\begin{figure}[H]
\centering
\begin{tikzpicture}[x=8.0cm,y=.53cm,>=Stealth,font=\small]
\draw[->] (-.03,0)--(1.22,0) node[right] {depth$/b$};
\foreach \x/\lab in {0/0,.3/{1-\ell},.7/{1-\tau},1/1}
 {\draw (\x,.07)--(\x,-.07) node[below] {$\lab$};}
\foreach \y/\name/\left in {1/V/0,2/U/.3,3/Y/.7,4/X/1}
 {\node[left] at (-.03,\y) {$\name$};
  \draw[gray,thick] (0,\y)--(\left,\y);
  \filldraw[fill=linkblue!22,draw=linkblue] (\left,\y-.19) rectangle +(0.12,.38);}
\draw[<->] (1,4.55)--(1.12,4.55) node[midway,above] {$\zeta$};
\node[anchor=west] at (0,5.55) {Distinct speeds: $0<\tau<\ell<1$};
\begin{scope}[yshift=-5.2cm]
\draw[->] (-.03,0)--(1.22,0) node[right] {depth$/b$};
\foreach \x/\lab in {0/0,.5/{1/2},1/1}
 {\draw (\x,.07)--(\x,-.07) node[below] {$\lab$};}
\foreach \y/\name/\left in {1/V/0,2/U/.5,3/Y/.5,4/X/1}
 {\node[left] at (-.03,\y) {$\name$};
  \draw[gray,thick] (0,\y)--(\left,\y);
  \filldraw[fill=linkblue!22,draw=linkblue] (\left,\y-.19) rectangle +(0.12,.38);}
\draw[dashed,linkblue] (.5,1.65) rectangle (.62,3.35);
\node[anchor=west] at (.66,2.5) {one joint middle group};
\node[anchor=west] at (0,5.05) {Tied speeds: $\tau=\ell=1/2$};
\end{scope}
\end{tikzpicture}
\caption{Schematic rebased coordinate windows for the local matrix
experiment in Section~\ref{sec:projections}. Gray segments show the base observations;
shaded windows show the additional refinements. Relative to the original offset, the base
of axis $i$ is shifted by $(1-k_i)b$, and its additional observation
has depth $\zeta b$. When $\zeta$ is smaller than every nonzero speed
separation, the windows of distinct groups are disjoint. At each window,
all groups with smaller $k_i$ are already observed in the full base
cell. The other distinct-speed order exchanges $Y$ and $U$.
Tied coordinates overlap and must retain their joint rate.}
\label{fig:entropy-layers}
\end{figure}


\begin{lemma}[Local counts and conditional Frostman laws]
\label{lem:blowup-counts}
There is a triangular choice $b_n\downarrow0$, $N_n\to\infty$,
with $D_n=N_n^{b_n}\to\infty$, and uniform refinements of the
single-event graphs, chosen independently of the queried offset,
with the following properties at almost every admissible offset
$a$. At stage $n$, evaluate the finite experiment at a prescribed
grid offset $a_n$ with $|a_n-a|=o(b_n)$; all rates below are
evaluated at the limiting point $(c,a)$. Conditional on
\[
 \mathcal B=(E_0,J_0,L),
\]
the entropy of $\mathcal O_{\mathbf u}$, divided by $\log D$, is
\begin{equation}
 \mathcal R(\mathbf u)+o(1).
\label{pr:local-entropy-formula}
\end{equation}
For distinct speeds,
\[
 \mathcal R(\mathbf u)=\sum_i j_i u_i.
\]
For tied middle speeds, put
\[
 r_Y=\phi_{\{X,Y\}}-\phi_{\{X\}},\qquad
 r_U=\phi_{\{X,U\}}-\phi_{\{X\}}.
\]
Then
\begin{equation}
 \mathcal R(\mathbf u)
   =a_1u_X+v_1u_V+
 \begin{cases}
    r_Y(u_Y-u_U)+m_1u_U,&u_Y\ge u_U,\\
    r_U(u_U-u_Y)+m_1u_Y,&u_U\ge u_Y.
 \end{cases}
\label{pr:tied-rectangle}
\end{equation}
These formulas hold for joint and unequal-depth observations.
Their occupied conditional cells have probabilities
\begin{equation}
 \Prob(\mathcal O_{\mathbf u}=o\mid\mathcal B)
      =D^{-\mathcal R(\mathbf u)+o(1)}
\label{pr:local-cell-masses}
\end{equation}
uniformly on the regularized graph at every tested depth.
Ratios of joint and conditioning cell probabilities give the
corresponding conditional formulas.

In particular, fix nuisance axes at relative depth $\zeta$ and
observe a selected pair to any depth $0\le u\le\zeta$.
For distinct speeds its conditional point law is Frostman of
exponent equal to the sum of the two selected rates, with subpower
constant, through resolution $D^{-\zeta}$.
For tied middle speeds, the additional $U$-rate after observing
$Y$ to depth $\zeta$ is $m_1-r_Y$; the pair $(X,V)$ after fixing
$Y,U$ has exponent $a_1+v_1$.

The same triangular choice supplies, conditionally on the original
index and $E_0$, the time and normal-parameter Frostman bounds of
exponents $s_0$ and $\gamma$ in relative units.
When $\tau=\ell$, the joint parameter square cells have masses
$D^{-s_0u+o(1)}$ for $0\le u\le\zeta_{\max}$.
\end{lemma}

\begin{proof}
We prove the rectangular entropy formula first. We then realize its
values as uniform finite cell masses and verify the time and normal
laws needed to choose a second event.

\paragraph{Rectangular entropy increments.}
Fix $\sigma>0$ and, temporarily, take $b$ so small that
$C_0b<\sigma$ and $c+\sigma<1$.
Condition additionally on $e_{c+\sigma}$ and read all forms in
that finer frame. Given $L$, Lemma~\ref{lem:pr-relative-frame}
shows that the required observations in the two frames determine
one another up to bounded overlaps. The full starting cell is
included in this comparison; errors from unobserved components
are below the finest tested width.

Apply the rectangular layering conclusion of
Lemma~\ref{lem:gap-density}. The starting offsets of the four
increment intervals are
\[
 a+(1-k_i)b.
\]
Intervals belonging to different speed groups are disjoint, because
$\zeta_{\max}<|k_i-k_j|$.
During the interval for a group of speed $k$, every group of smaller
speed has already been required to a depth beyond that interval,
even in its \emph{base} observation. Every group of larger speed
ends below the interval. If $S$ is the union of the groups of smaller
speed and $G$ is the current group, subtraction of the base rectangle from the refined
rectangle leaves the density
$\phi_{S\cup G}^{\sigma}-\phi_S^{\sigma}=:j_G^{\sigma}$.
Thus the group's contribution is its ordered conditional rate times
its refinement depth.
This reasoning applies both to the order $X,Y,U,V$, when
$\tau<\ell$, and to $X,U,Y,V$, when $\ell<\tau$.
It gives \eqref{pr:local-entropy-formula} with the
$\sigma$-conditioned rates, to error $o(b)$ in entropy divided by
$\log N$.

For a tied group, first reveal $Y$ and $U$ together to the smaller
of their two depths; then refine only the deeper coordinate. The
common part costs $m_1$ times that depth. The remaining
increment uses $r_Y$ or $r_U$, according to the deeper coordinate.
This gives \eqref{pr:tied-rectangle}, again first with
$\sigma$-conditioned rates.
Equivalently, one may subtract the layering formulas for a joint
rectangle and its conditioning rectangle. All their thresholds lie
within $a\pm Cb$ for a fixed $C$.

We next remove the extra label. Given $L$, the joint and
conditioning observations in question are sandwiched between full
$e_c$-frame observations at offsets $a-Cb$ and $a+Cb$, up to bounded
overlaps. By Lemma~\ref{lem:remove-label}, the change of a
conditional increment upon adding $e_{c+\sigma}$ is bounded by
the incremental label information on this full-cell interval.
For almost every $a$, its upper first-order density tends to zero
as $\sigma\downarrow0$.
Consequently the limsup of the error divided by $b$, as $b\to0$
with $\sigma$ fixed, is bounded by a constant times that density.
Now let $\sigma\downarrow0$ through a countable sequence.
The rates converge to $\phi_S$ and $j_G$, and the error is $o(b)$.
This proves the asserted formulas for the limiting normalized
entropies without requiring a convergence speed uniform in $a$.

\paragraph{Uniform finite cell masses.}
Here is a simultaneous implementation on finite configurations.
Fix a summable sequence of positive rational $b_n\downarrow0$.
At stage $n$, prescribe a finite offset grid of spacing $o(b_n)$
on the compact interval under consideration, and a finite grid of
relative depths whose union is dense in
$[0,\zeta_{\max}]$. Include the following in the countable family
of tests: $E_0,J_0,L$; every subset of the axis observations at
equal and unequal depths; their joint and conditioning cells;
the full-cell sandwiches just used; and the same observations with
the fixed additional labels $e_{c+\sigma}$.
Include also the within-index parameter-cell tests described below.

Use Lemma~\ref{lem:regularization} on finitely many tests at a time,
including the deficient-cell deletions relative to their original
mass. This makes the joint and conditional cell masses uniform in
exponent. Previously uniform tests keep their limits: surviving
cells retain their original masses within the prescribed subpower
factor, and the total restriction has the same order of loss.
Thus the already defined stable entropy function and its rates do
not change.
Take subsequential limits on this countable family first.
Their values satisfy the estimates established above with fixed
$\sigma$ and then $\sigma\downarrow0$.

Finally choose $N_n$ sufficiently large along the resulting
configurations that the first $n$ finite collections of tests have
errors $o(b_n)$ in exponents normalized by $\log N_n$.
The finite pigeonholings and cleaning can have total loss
$N_n^{-o(b_n)}$: for a fixed finite number of tests their number of
mass levels grows at most polynomially in $\log N_n$ and in the
inverse requested accuracy, so increasing $N_n$ makes their
normalized logarithmic cost smaller than the chosen multiple of
$b_n$. Choose the deletion cutoffs below that retention while still
with exponent $o(b_n)$.
Also require $b_n\log N_n\to\infty$.
For each generic $a$, use its nearest offset-grid point.
The stable full-cell and finer-label comparisons change the
exponents by $o(b_n)$ at this replacement.
The first-order density estimates hold at the actual generic point
for every sequence $b_n\to0$; no uniform differentiability rate
over all such points was used.

Uniform joint cell masses now prove
\eqref{pr:local-cell-masses}, rather than only its entropy average.
For example, in the tied case and $u\le\zeta$,
\begin{align*}
 &\mathrm H( U_u,Y_\zeta\mid\mathcal B)
       -\mathrm H(Y_\zeta\mid\mathcal B)\\
 &\hspace{25mm}=(m_1-r_Y)u\log D+o(\log D).
\end{align*}
The probability of an occupied $U_u$ cell within an occupied
$Y_\zeta$ cell is the ratio of their joint and marginal
probabilities, hence $D^{-(m_1-r_Y)u+o(1)}$.
For distinct speeds the same subtraction removes all nuisance
terms and leaves $u$ times the sum of the selected rates.
For $(X,V)$ with $Y,U$ fixed, the tied-group contribution cancels
and leaves $(a_1+v_1)u$.
These statements hold at all the tested intermediate depths.
Bounded square coverings and increasingly fine depth meshes give
the asserted Frostman upper bounds at every intermediate radius.

\paragraph{Time and normal parameters.}
It remains to verify the parameter hypotheses for these finite
experiments. Lemma~\ref{lem:angular-frostman} and the uniform time
branching give the two upper bounds conditional on the index and
$E_0$. Nondeficient $E_0$ cells have the same exponent as their
time blocks, since each index uses only subpower many angular bins
there. Add these conditional cells and their refinements to the
tests before choosing the triangular sequence.
When $\tau=\ell$, refining the label from $c+b$ to
$c+b+ub/\ell$ multiplies both relative widths by $D^{-u}$.
The time mass is $D^{-s_0u+o(1)}$. Within each of its time blocks,
the normal range occupies only subpower many corresponding
angular bins. Its nondeficient occupied joint cells therefore have
that same mass exponent. Cleaning them gives the lower as well as
the upper bound, uniformly through the required depths.
All of this cleaning precedes the choice of a pair law; a later
bias is never applied to an uncontrolled exceptional set of
single-event cells.
\end{proof}

\subsection{The product-slope pinning input}

When $\tau=\ell$, the local count formula supplies the joint middle
rate $m_1$, rather than two separate ordered middle rates. To bound
$v_1$ in this case, freeze both $Y$ and $U$ in the first formula of
\eqref{pr:relative-frame}. The remaining projection is
\[
 X'=X+(\Delta t\Delta\vartheta)V
                    +\Delta t\,U+\Delta\vartheta\,Y.
\]
The last two terms are a translation once the nuisance coordinates
and the two labels are fixed. The slope is the product of two
parameter differences; their separate Frostman bounds do not assert
that this product is nonconcentrated. The following lemma supplies
that bound while controlling the change in each single-event law.

In its statement $D\to\infty$ is a resolution parameter and $\zeta$
is a positive depth. Neither denotes a critical parameter.
A probability $\mu$ on a bounded square is called uniformly
$s$-regular through depth $\zeta_{\max}$ if, for every occupied
nested dyadic square $Q$ of side comparable to $D^{-u}$,
\begin{equation}
 D^{-su-\epsilon_D}\le\mu(Q)\le D^{-su+\epsilon_D},
 \qquad 0\le u\le\zeta_{\max},\qquad \epsilon_D\longrightarrow0.
\label{pr:square-regularity}
\end{equation}
Fixed factors from dyadic rounding can be included in
$D^{\epsilon_D}$. It is equivalent here to impose these estimates
on depth meshes whose spacing tends to zero: monotonicity gives
\eqref{pr:square-regularity}, with a slightly larger error.
Uniformity for an ensemble means that the same error works for all
its members after deleting a set of initial indices of probability
tending to zero.

\begin{lemma}[Pinning with controlled marginals]
\label{lem:pinning}
Let $0<s\le1$, $\gamma'>0$, and $0<\zeta_{\max}\le1$ be fixed.
For each $D$, let $I$ have a finite probability distribution $\pi_D$,
and let $\mu_i$ be finitely supported probabilities on a fixed
bounded square with coordinates $(t,\vartheta)$.
Suppose uniformly in $i$ that \eqref{pr:square-regularity} holds and
that the two coordinate marginals satisfy
\begin{equation}
 \mu_i\{t\in J_r\}\le D^{o(1)}r^s,
 \qquad
 \mu_i\{\vartheta\in J_r\}\le D^{o(1)}r^{\gamma'}
 \quad(D^{-\zeta_{\max}}\le r\le1)
\label{pr:coordinate-frostman}
\end{equation}
for all intervals $J_r$ of radius $r$.
For every $\epsilon>0$, there are a subsequence, a fixed rational
$\zeta\in(0,\zeta_{\max}]$, and a probability $\nu_D$ on triples
$(i,\xi,\upsilon)$, with $\xi,\upsilon\in\supp\mu_i$, such that:
\begin{enumerate}[label=\textup{(\roman*)}]
\item each marginal $(i,\xi)$ and $(i,\upsilon)$ is dominated by
$D^{\epsilon\zeta}\pi_D(di)\mu_i$;
\item for $\nu_D$-almost every $(i,\xi)$, the conditional law of
\[
 \mathfrak p_\xi(\upsilon)
   =(t_\upsilon-t_\xi)(\vartheta_\upsilon-\vartheta_\xi)
\]
satisfies
\begin{equation}
 \nu_D\{\mathfrak p_\xi(\upsilon)\in J_r\mid i,\xi\}
 \le C D^{\epsilon\zeta}r^{s-\epsilon},
 \qquad D^{-\zeta}\le r\le1.
\label{pr:pinned-frostman}
\end{equation}
\end{enumerate}
The pair law can be chosen using only the initial index and its
parameter measure. In particular it need not depend on any later
matrix-cell or offset test.
\end{lemma}

The proof of Lemma~\ref{lem:pinning} is in
Section~\ref{sec:product-pinning-proof}.

\subsection{Two events on one index}

Let $\mu$ be the probability on the regularized single-event graph.
First draw $(I,E_0)$ from its marginal under $\mu$. For the
\emph{ordinary pair law}, draw $\xi_0,\xi_1$ conditionally
independently from the original law given $(I,E_0)$. Both events use
the matrix attached to $I$. Their conditional independence does not
make that matrix independent of their frame labels.
Let $L_0,L_1$ be their labels at the precision in
\eqref{pr:fine-label}.
Conditional on $(I,E_0,\xi_0)$, the difference of the second time
from the first is $s_0$-Frostman, and the corresponding normal
difference is $\gamma$-Frostman, in relative units. Translation
does not change either estimate. Reading the differences from
$L_0,L_1$ changes them by less than the testing resolution, by
Lemma~\ref{lem:pr-relative-frame}. Bounded enlargement of intervals
therefore preserves their Frostman bounds.

In the tied case we need one additional coupling. Apply
Lemma~\ref{lem:pinning} to the within-index parameter laws given
$(I,E_0)$, with $s=s_0$, using the last assertion of
Lemma~\ref{lem:blowup-counts}. For any fixed sufficiently small
$\epsilon>0$, it gives, along a further subsequence, a fixed
rational depth $\zeta>0$. Both single-event marginals are bounded
by $D^{\epsilon\zeta}$ times their original laws. Conditional on
$(I,E_0,\xi_0)$ the product slope
\begin{equation}
 \Theta_*=(t_1-t_0)(\vartheta_1-\vartheta_0)
\label{pr:product-slope}
\end{equation}
is $(s_0-\epsilon)$-Frostman with constant
$CD^{\epsilon\zeta}$ through resolution $D^{-\zeta}$.
All coordinates here are in the relative parameter units of
$E_0$. The product is Lipschitz on this fixed bounded square, so
the fine-label replacement preserves this estimate as well.

The original graph and these pair distributions are chosen before
querying an offset. We can prescribe countably many loss
tolerances and projection tests. Each can use a further subsequence;
intersecting their resulting full-measure sets of offsets still
gives a full-measure set. No one choice of pair distribution is
asserted to work for every tolerance.

The same-index relation forces the matrix, read in the second frame,
to belong to the support of that frame's original coordinate law. We
will compare this certain event with its probability under the product
of the conditional matrix and label marginals. A small probability
under that product would require positive mutual information. The
next estimate shows that the short matrix-refinement gap contains
asymptotically less information than such a comparison would require.
It applies equally to the ordinary and pinned pair laws. The finite
pair law may depend on $b,N,I,E_0$, but it must be fixed while the offset
is integrated.

\begin{lemma}[Pair information on the refinement gap]
\label{lem:pair-information}
Let $K=1+\zeta_{\max}$, and let $J_f$ be the full observation in
the $E_0$-frame obtained by increasing every threshold of $J_0$
by $Kb$. Along a triangular subsequence as in
Lemma~\ref{lem:blowup-counts}, for almost every admissible offset,
\begin{equation}
 \mathrm I(J_f;L_0,L_1\mid E_0,J_0)=o(\log D).
\label{pr:pair-information}
\end{equation}
Here $J_0,J_f$ are evaluated at the nearest prescribed offset-grid
point. The assertion holds for every one of the countably many
pair experiments just described, with ordinary Shannon units.
\end{lemma}

\begin{proof}
Given $E_0$, each label $L_i$ has at most
$C N^{C_1b}$ possible values, where $C_1$ depends on $C_0$ and
the fixed speeds. This is an alphabet bound from the time and
normal grids, independent of the probabilities of these labels.
It remains true under either marginal bias. Consequently
\begin{equation}
 \mathrm H(L_0,L_1\mid E_0)\le C_2(1+b\log N).
\label{pr:label-alphabet}
\end{equation}

Temporarily write $J_0(a)$ for the rebase grid at offset $a$ and
set
\[
 f(a)=\mathrm I(J_0(a);L_0,L_1\mid E_0).
\]
This is nondecreasing, because the matrix grids are nested as $a$
increases. Its range is contained in
$[0,C_2(1+b\log N)]$. The chain rule gives the exact identity
\begin{equation}
 f(a+Kb)-f(a)
 =\mathrm I(J_0(a+Kb);L_0,L_1\mid E_0,J_0(a)).
\label{pr:information-increment}
\end{equation}
For a nondecreasing function with range of length $H_*$,
\[
 \int_A^B[f(a+q)-f(a)]\,da\le qH_*
 \qquad(q>0),
\]
as follows either by changing variables in the two integrals or
by integrating its nonnegative increment measure.
For a grid replacement $a_n$ with $|a_n-a|=o(b_n)$, monotonicity
also bounds the left side of \eqref{pr:information-increment}
by
\[
 f(a+(K+1)b)-f(a-b)
\]
for large $n$. Integrate on a compact interval whose slightly
larger neighborhood stays admissible. Equations
\eqref{pr:label-alphabet} and \eqref{pr:information-increment}
give
\begin{equation}
 \int
 \frac{\mathrm I(J_f;L_0,L_1\mid E_0,J_0)}{\log D}\,da
 \le C_3\left(b+\frac1{\log N}\right).
\label{pr:information-integral}
\end{equation}
The left side uses the stepwise chosen grid offsets. The
containing-gap estimate proves the bound regardless of the sizes
of the grid's inverse images.

Choose $b_n$ summable, and increase $N_n$ so that
$1/\log N_n$ is summable as well as satisfying all the finite
tests of Lemma~\ref{lem:blowup-counts}. Tonelli's theorem applied
to the nonnegative functions in
\eqref{pr:information-integral} shows that their sum is finite
almost everywhere. In particular the summands tend to zero,
which proves \eqref{pr:pair-information}.
A further subsequence, such as the one used for pinning, preserves
this conclusion. For countably many pair experiments use the same
integral argument for each and intersect the full-measure sets.
Finally exhaust the admissible offset domain by compact intervals.
\end{proof}

We next record explicitly how the marginal bias is used. Denote
the first single-event marginal of a chosen pair law by $\nu$.
Suppose
\begin{equation}
 \nu\le D^{\epsilon\zeta}\mu.
\label{pr:marginal-domination}
\end{equation}
Let $\mathcal B_0=(E_0,J_0,L_0)$.
Discard base values $B$ for which
$\nu(B)<D^{-\epsilon\zeta}\mu(B)$.
Their total $\nu$-probability is at most $D^{-\epsilon\zeta}$.
For every other base, division of
\eqref{pr:marginal-domination} by its base mass gives
\begin{equation}
 \nu(\,\cdot\mid B)
       \le D^{2\epsilon\zeta}\mu(\,\cdot\mid B).
\label{pr:conditional-domination}
\end{equation}
This also holds for the conditional matrix marginals.
For the ordinary pair law the first marginal equals
$\mu$, and no base deletion is needed.

\subsection{Comparing the coupling with a product law}

For a fixed base $B=(E_0,J_0,L_0)$, the actual pair experiment
defines a joint law of the matrix and $L_1$.
The conditional slope law is obtained by mixing its laws given
$(I,E_0,\xi_0)$. Since $J_0$ is determined by the matrix and
$E_0$, this conditioning and mixture preserve each of the uniform
time, angular, or pinned-product Frostman bounds above.
We compare the joint law with the product of its two conditional
marginals. Independence is imposed only in this comparison law.

The planar estimate we use keeps the entire direction label: two
labels with the same slope may prescribe different target intervals.
Its point and label measures are independent in the product below.
The proof is given in Section~\ref{sec:robust-projection-proof}.

\begin{lemma}[Robust planar projections]\label{lem:robust-projection}
Fix
\[
 0<s\le1,\qquad 0<t\le2,\qquad
 0\le u<h_{\mathrm{pr}}:=\min\{t,(s+t)/2,1\}.
\]
There is $\kappa_0=\kappa_0(s,t,u)>0$ with the following property.
For every $R,C_0\ge1$ and $0<\kappa\le\kappa_0$, there is
$\delta_0>0$ such that the following holds for
$0<\delta<\delta_0$. Let $\mu$ be a Borel subprobability measure
on $[-R,R]^2$, let $\nu$ be a subprobability measure on an
arbitrary finite label set $\Lambda$, and let
$\vartheta:\Lambda\to[-R,R]$. Suppose
\begin{equation}\label{eq:robust-projection-frostman}
 \mu(B(p,r))\le\delta^{-\kappa}r^t,
 \qquad
 (\vartheta_*\nu)(B(a,r))\le\delta^{-\kappa}r^s
 \quad(\delta\le r\le1).
\end{equation}
For each $\lambda\in\Lambda$, let $A_\lambda$ be a union of at
most $C_0\delta^{-u}$ intervals of length at most $C_0\delta$.
Then
\begin{equation}\label{eq:robust-projection-product}
 \sum_{\lambda\in\Lambda}\nu(\{\lambda\})
 \mu\{(x,y):x+\vartheta(\lambda)y\in A_\lambda\}
 <\delta^{2\kappa}.
\end{equation}
Consequently there cannot be a set $B\subset\Lambda$ with
$\nu(B)\ge\delta^\kappa$ such that for every $\lambda\in B$
there is a $\mu$-measurable set $E_\lambda$ with
\begin{equation}\label{eq:robust-projection-subset}
 \mu(E_\lambda)\ge\delta^\kappa,
 \qquad
 \{x+\vartheta(\lambda)y:(x,y)\in E_\lambda\}
 \subset A_\lambda.
\end{equation}
\end{lemma}

\begin{lemma}[Transfer of a projection bound to rates]
\label{lem:pr-projection-transfer}
Fix the original single-event law $\mu$ and one of the same-index
pair laws above, with the conditional marginal domination already
established. Suppose that after freezing
the nuisance axes at relative depth $\zeta$, the original
conditional law of two source coordinates is Frostman with
exponent $t>0$ and subpower constant through resolution
$\delta=D^{-\zeta}$, uniformly in its occupied nuisance fibers.
Suppose that, in the pair law conditional on $B$, the changing
projection slope has an available positive Frostman exponent $s\le1$.
In a pinned experiment, ``available'' means exponents tending to
$s$ with arbitrarily small fixed marginal and Frostman losses.
Suppose further that the target coordinate in the $L_1$-frame
has, conditional on $(E_0,J_0,L_1)$, at most
$D^{a_*\zeta+o(1)}$ occupied bins of width comparable to
$D^{-\zeta}$.
Then
\begin{equation}
 a_*\ge \min\{t,(t+s)/2,1\}.
\label{pr:projection-transfer}
\end{equation}
\end{lemma}

\begin{proof}
Suppose the inequality fails strictly. Choose fixed
$0<t_*<t$, $0<s_*<s$, and $u>a_*$, still satisfying
\begin{equation}
 u<\min\{t_*,(t_*+s_*)/2,1\}.
\label{pr:strict-projection-margin}
\end{equation}
The choices are possible by continuity, and are made before any
pinning loss is chosen. Let $\kappa>0$ be an allowance from
Lemma~\ref{lem:robust-projection} for these three fixed exponents.
In the pinned case choose the loss tolerance $\epsilon$ so small
that
\begin{equation}
 \epsilon<s-s_*,\qquad 2\epsilon<\kappa/2.
\label{pr:loss-order}
\end{equation}
Choose the pinned coupling and its fixed positive depth $\zeta$
with this tolerance. For the ordinary pair law any fixed rational
$0<\zeta\le\zeta_{\max}$ works. The depth does not subsequently
shrink with $D$. Subpower errors in the original single-event
estimates therefore tend to zero also after division by $\zeta$.

\paragraph{The product probability bound.}
Fix a good base $B$, and initially use the \emph{original}
conditional matrix law $\mu(\,\cdot\mid B)$, independently of
the pair experiment's conditional law of $L_1$.
Within this original matrix law, freeze the nuisance bins at
depth $\zeta$. Lemma~\ref{lem:blowup-counts} gives the point
Frostman bounds, with exponent $t_*$, in every occupied fiber.
For sufficiently large $D$ their constants are smaller than
$\delta^{-\kappa}$. The slope bounds, with exponent $s_*$, have
the same allowed constant: in the pinned case they cost
$C\delta^{-\epsilon}$, and for the ordinary pair law a subpower.

The frame formula \eqref{pr:relative-frame} writes the target
coordinate as a two-coordinate scalar projection plus a
translation depending on the frozen nuisance bins and the full
label $L_1$. The error from a nuisance bin has size $O(\delta)$,
since all coefficients are bounded. For each $L_1$, enlarge all
its allowable target bins by this error and subtract the known
translation. Their number is at most $C\delta^{-u}$, because
$u>a_*$. Their lengths are at most $C\delta$.
The label may contain more information than the projection slope;
Lemma~\ref{lem:robust-projection} permits precisely this dependence
of the target intervals on the entire label.

Apply that lemma in each original nuisance fiber, with fixed
bounded coordinate charts. Its independent product hit
probability is less than $\delta^{2\kappa}$, uniformly in the
fiber. Integrating over the original nuisance distribution gives
the same bound for the entire original conditional matrix law.
Only now replace that whole matrix law by its biased marginal.
Equation~\eqref{pr:conditional-domination} bounds the new product
hit probability by
\begin{equation}
 \delta^{-2\epsilon}\delta^{2\kappa}
       \le\delta^{\kappa}
\label{pr:biased-product-hit}
\end{equation}
for large $D$. For the ordinary pair law the domination factor is
one. This argument does not assume that biased nuisance fibers
have the original Frostman estimates, or that a slope remains
independent after conditioning inside the actual coupling.

\paragraph{The support event for the same-index law.}
Under the true coupling, the target matrix coordinate always
belongs to the allowed support for $(E_0,J_0,L_1)$.
Both pair marginals are supported on the original single-event
graph and both events have exactly the same index matrix. Thus
this support test has probability one in the true coupling.
It requires no probability lower bound for target cells after
the change of measure.

To make it a finite information test, replace the matrix by a
representative of its $J_f$ cell. The relative coordinate error
is $O(D^{-(1+\zeta_{\max})})$, which is smaller than
$D^{-\zeta}$ by a fixed positive power. The bounded frame
coefficients and the precision of $L_0,L_1$ leave an
$O(\delta)$ error. Enlarge the target bins once more.
The resulting event $A_B$ depends only on $(J_f,L_1)$ at fixed
$B$, has conditional probability one in the true coupling, and
still has probability at most $\delta^\kappa$ under the product
of the conditional marginals of $J_f$ and $L_1$.
For the last assertion, sample the exact matrix from its
conditional marginal and then round it. A rounded hit implies
an exact hit in one further bounded enlargement, already allowed
in the preceding application of Lemma~\ref{lem:robust-projection}.

\paragraph{The information contradiction.}
Let $P_B$ be the conditional joint law of $(J_f,L_1)$ and $Q_B$
the product of its marginals. Data processing of relative entropy
by the indicator of $A_B$ gives
\[
 \mathrm I(J_f;L_1\mid B)
   =\mathrm D_{\mathrm{KL}}(P_B\Vert Q_B)
   \ge \log\frac1{Q_B(A_B)}
   \ge \kappa\zeta\log D.
\]
The good bases have pair probability tending to one by
\eqref{pr:conditional-domination}. Averaging the last inequality
therefore gives
\begin{equation}
 \mathrm I(J_f;L_1\mid E_0,J_0,L_0)
   \ge (1-o(1))\kappa\zeta\log D.
\label{pr:information-contradiction}
\end{equation}
On the other hand the chain rule and nonnegativity imply
\[
 \mathrm I(J_f;L_1\mid E_0,J_0,L_0)
 \le \mathrm I(J_f;L_0,L_1\mid E_0,J_0)
 =o(\log D)
\]
by Lemma~\ref{lem:pair-information}. Since $\kappa\zeta>0$
is fixed, this contradicts
\eqref{pr:information-contradiction} and proves the lemma.
\end{proof}

\subsection{Deriving the individual constraints}

\begin{proof}[Proof of Proposition~\ref{prop:projection-rates}]
Fix a good interior $c$. Work at offsets where the entropy
densities, the local counts, and the pair information conclusions
all hold. Strict violations of any inequality can be witnessed
by rational choices of the smaller exponents and loss tolerances
in \eqref{pr:strict-projection-margin}--\eqref{pr:loss-order}.
There are countably many such choices, so the preceding
constructions apply on a common full-measure set of offsets.
Fubini then gives the asserted almost-everywhere statement in
$(c,a)$.

For distinct speeds the four tests needed below are displayed in
Table~\ref{tab:pr-tests}. Each uses the ordinary pair law. Source
exponents are conditional on the listed
nuisance bins in the original matrix law. They follow from the
all-depth probability ratios in Lemma~\ref{lem:blowup-counts},
in either of the two possible speed orders.

\begin{table}[htbp]
\centering
\small
\begin{tabular}{@{}llllll@{}}
\toprule
Target & Source pair & Frozen axes & Slope & Source exponent
 & Target exponent\\
\midrule
$X$ & $(X,Y)$ & $U,V$ & $\Delta\vartheta$ & $a_1+y_1$ & $a_1$\\
$X$ & $(X,U)$ & $Y,V$ & $\Delta t$ & $a_1+u_1$ & $a_1$\\
$U$ & $(U,V)$ & $X,Y$ & $\Delta\vartheta$ & $u_1+v_1$ & $u_1$\\
$Y$ & $(Y,V)$ & $X,U$ & $\Delta t$ & $y_1+v_1$ & $y_1$\\
\bottomrule
\end{tabular}
\caption{The four ordinary-pair tests for distinct speeds.
The angular and time slope exponents are respectively
$\gamma$ and $s_0$.}
\label{tab:pr-tests}
\end{table}

For example the first formula of
\eqref{pr:relative-frame}, with $U,V$ frozen, is
\[
 X'=X+\Delta\vartheta Y
       +\Delta t U+\Delta t\Delta\vartheta V.
\]
The last two terms are a full-label-dependent translation up to
$O(D^{-\zeta})$. With $Y,V$ frozen, the same expression is
$X+\Delta t U$ plus such a translation. The second and third
formulas give the other two rows. Thus every row has exactly the
form required in Lemma~\ref{lem:pr-projection-transfer}.

It is useful to spell out the numerical implication. For
$0\le a,b\le1$ and $s>0$, suppose the test with source exponent
$a+b>0$ gives
\begin{equation}
 a\ge\min\{a+b,(a+b+s)/2,1\}.
\label{pr:elementary-test}
\end{equation}
If $a<1$ and $b>0$, both the first and the third entries of this
minimum exceed $a$. The second must therefore be at most $a$,
so $b\le a-s$. If $b=0$ there is no further restriction.
In either case
\begin{equation}
 a<1\quad\Longrightarrow\quad b\le(a-s)_+.
\label{pr:elementary-consequence}
\end{equation}
When $a+b=0$, that conclusion is immediate without applying the
projection lemma.

The first row gives $y_1\le(a_1-\gamma)_+$ when $a_1<1$.
Together with $y_1\le1$, this implies $a_1\ge y_1$ also when
$a_1=1$. The second row gives
$u_1\le(a_1-s_0)_+$ whenever $a_1<1$.
The third row analogously gives $u_1\ge v_1$.
For the last row, if $v_1>0$ and $y_1<1$,
\eqref{pr:elementary-consequence} gives $y_1\ge s_0+v_1$.
If $y_1=1$, the required
$y_1\ge\min\{1,s_0+v_1\}$ holds automatically.
These are all assertions of \eqref{pr:distinct-constraints}.

Suppose now that $\tau=\ell$. Put
\[
 r_Y=\phi_{\{X,Y\}}-\phi_{\{X\}}.
\]
Conditional entropy monotonicity and the one-dimensional
alphabet bound give
$0\le r_Y\le1$ and $0\le m_1-r_Y\le1$.
Use the second row's time projection with $Y,V$ frozen.
For every $0\le u\le\zeta$, the tied rectangle formula gives
the conditional source exponent
\[
 t=a_1+(m_1-r_Y).
\]
The target exponent is $a_1$ and the slope exponent is $s_0$.
Equation~\eqref{pr:elementary-consequence}, with
$b=m_1-r_Y$, yields
\[
 a_1<1\quad\Longrightarrow\quad
 m_1-r_Y\le(a_1-s_0)_+.
\]
Since $r_Y\le1$, this proves the bound for $m_1$ in
\eqref{pr:tied-constraints}.

Finally freeze $Y,U$ and use the pinned pair distribution.
The first formula in \eqref{pr:relative-frame} becomes
\[
 X'=X+(\Delta t\Delta\vartheta)V
                 +\Delta t U+\Delta\vartheta Y.
\]
The source pair is $(X,V)$, of conditional exponent
$a_1+v_1$, by Lemma~\ref{lem:blowup-counts}.
Its target exponent is $a_1$.
The product slope has available exponent $s_0$ by
Lemma~\ref{lem:pinning} and the preceding coupling construction.
Lemma~\ref{lem:pr-projection-transfer} includes both its marginal
bias and its small direction loss. Thus
\eqref{pr:elementary-consequence} gives
$v_1\le(a_1-s_0)_+$ when $a_1<1$, as required.

Every construction above takes place on a compact subset of
$\mathcal U$ from \eqref{eq:entropy-domain}.
Indeed the smallest and largest of the four unshifted thresholds
are respectively
\[
 (1-w)(1-c)+a,\qquad 1-w(1-c)+a.
\]
If $w<1$, a compact interior $c$-interval admits sufficiently
small offsets of either sign, and the $O(b)$ rebase and
refinement shifts stay in range. If $w=1$, only positive small
offsets are interior. This proves the stated conclusion for every
permitted sign and completes the proposition.
\end{proof}
